\honest{When None Dare Urge Restraint}{Eliezer Yudkowsky}{2007}

On the morning of September 11, my email client showed me the news. My first thoughts were that I really was living in the Future, and ``Thank goodness it wasn't nuclear.'' My third was: ``The overreaction to this will be ten times worse than the original event.'' That was ``a vast understatement''; it is hard to aim pessimism low enough to be pleasantly surprised as often as unpleasantly. I do not say what I am measuring. \nb{By American military deaths in Iraq, 4,000 by March 2008, the ratio to the 2,977 victims of 11 September is under two. By CNN's estimates of Iraqi deaths in the same month, ``about 80,000 to the hundreds of thousands,'' it is far over ten.}

I knew at once that no one ``would dare to be the voice of restraint, of proportionate response.'' At first six thousand were thought dead. A politician who said that was an eighth of a year's American traffic deaths would have been asked to resign ``the same hour.'' \nb{On 14 September Barbara Lee told the House that ``some of us must urge the use of restraint,'' ``so that this does not spiral out of control.'' The resolution passed 420 to 1. Lee was called a ``traitor'' and received death threats, so the prediction held except for her.}

If everyone gains by stressing the pain and no one dares urge restraint, the response will overshoot ``whatever the appropriate level may be.'' This is the spiral of hate: whoever attacks the Enemy is a patriot, and whoever questions a negative claim about the Enemy is a traitor. But most complex statements are false, and so are most bad things you could say about anyone, even the worst person in the world.

My example is ``the suicide hijackers were cowards.'' It takes some courage to fly a plane into a building; of all their sins, cowardice was not one. Would I earn more credit by accusing al-Qaeda of killing Kennedy? \nb{President Bush had called the attacks ``cowardly acts'' that same day. Bill Maher and Susan Sontag said in public, within two weeks, that the hijackers were not cowards, and both paid for it. The post does not mention them.} Accuracy matters about anyone, Hitler included.

A defense with thousands of aircraft and hundreds of thousands of soldiers can do more damage than nineteen men with four airliners. The US spent ``billions of dollars and thousands of soldiers' lives shooting off its own foot.'' I do not name the war. Ignoring the attack ``would have been better than the real course of history.'' I do not say what an ignored attack would have led to. Politicians walked ``straight into al-Qaeda's trap.'' \nb{In 2004 bin Laden boasted that it had been easy ``to provoke and bait this administration.''}

At first there were smarter responses than I expected. A member of Congress, whose name I forget, said on camera that the first purpose of government is not the economy or health care but defending the country; that was not an applause light. Within two days, concern for image took over, and once restraint becomes unspeakable, fury and folly ``can only rise with time.'' \nb{In Gallup's poll a week before the post, 57\% of Americans said sending troops to Iraq had been a mistake, against 23\% in March 2003. By then restraint was no longer unspeakable, and the post does not say how that came about.}
